\documentclass[twocolumn]{article}
\usepackage[utf8]{inputenc}
\usepackage[turkish]{babel}
\usepackage{amsmath}
\usepackage{graphicx}
\usepackage{cite}
\usepackage{hyperref}
\usepackage{booktabs}

\title{NPU Mimarilerinde Graf Optimizasyonlarinin Ampirik Calisma Zamanina Etkisi: Operator Fuzyonu Uzerine Bir Zaman ve I/O Karmasikligi Analizi}
\author{Yusuf Talha ARABACI \\ Karabuk Universitesi, Yazilim Muhendisligi YL}
\date{Nisan 2026}

\begin{document}

\maketitle

\begin{abstract}
Modern yapay zeka hizlandiricilarinda (NPU) performans, hesaplama kapasitesinden ziyade bellek bant genisligi ile sinirlanmaktadir. Bu calismada, derin ogrenme modellerinin temsil edildigi Yonlu Asiklik Graflar (DAG) uzerinde uygulanan Operator Fuzyonu (Operator Fusion) isleminin ampirik etkileri incelenmistir. Deneysel sonuclar, dugum birlestirme isleminin teorik hesaplama karmasikligini degistirmemesine ragmen, I/O trafigini azaltarak sabit carpani ($c$) kuculttugunu gostermektedir.
\end{abstract}

\section{Giris}
Geleneksel algoritma analizinde kullanilan RAM modeli, bellek erisim maliyetlerini $\mathcal{O}(1)$ kabul eder. Ancak NPU donanimlarinda Bellek Duvari (Memory Wall) darbohazi nedeniyle ana bellek (DRAM) erisimi, hesaplama birimlerinden cok daha yavastir. Bu durum, algoritmalarin asimptotik analizinde I/O karmasikliginin onemini artirmaktadir.

\section{Teorik Analiz: DAG ve Dugum Birlestirme}
Bir sinir agi, matematiksel operatorlerin dugum ($V$) ve veri akisinin ayrit ($E$) oldugu bir $G=(V,E)$ DAG yapisidir.

\subsection{Dugum Birlestirme (Node Merging)}
Operator fuzyonu, algoritmik olarak ardisik dugumlerin birlestirilerek tek bir hesaplama cekirdegine (kernel) donusturulmesidir:
\begin{equation}
    f_{\text{fused}}(x) = f_2(f_1(x))
\end{equation}
Graf seviyesinde $|V|$ azalsa da toplam operasyon miktari korunur. Bu nedenle hesaplama karmasikligi sinifi yaklasik olarak $\mathcal{O}(N+E)$ kalir. Buna karsin, ara sonuc yaz/oku adimlari azalir ve I/O sabitleri kuculur.

\section{Metodoloji}
Calismada ONNX Runtime kullanilarak model bazli ampirik testler yurutmektedir. Iki senaryo karsilastirilmistir:
\begin{itemize}
    \item \textbf{Baseline:} \texttt{ORT\_DISABLE\_ALL} (fuzyon yok)
    \item \textbf{Optimized:} \texttt{ORT\_ENABLE\_EXTENDED} (fuzyon aktif)
\end{itemize}

Her senaryoda 100 iterasyon olcum yapilmis, operator seviyesinde profiling JSON dosyalari disariya alinmis ve operator bazli sure kirilimlari CSV olarak raporlanmistir.

\subsection{Deney Tasarimi Tablosu}
\begin{table}[h]
\centering
\begin{tabular}{|l|c|c|r|}
\hline
\textbf{Metrik} & \textbf{Baseline} & \textbf{Optimized} & \textbf{Degisim} \\ \hline
Dugum Sayisi ($|V|$) & 142 & 78 & -\%45 \\ \hline
Bellek Trafigi (GB/s) & 12.4 & 6.1 & -\%50 \\ \hline
Gecikme (ms) & 15.2 & 9.4 & -\%38 \\ \hline
\end{tabular}
\caption{Graf Optimizasyonunun Ampirik Sonuclari}
\end{table}

\section{Roofline Model ile Yorum}
Roofline modeline gore performans siniri su sekilde verilir:
\begin{equation}
P \leq \min(P_{\text{peak}},\ I \times B_{\text{peak}})
\end{equation}
Burada $I$ aritmetik yogunluk (FLOP/Byte), $B_{\text{peak}}$ bellek bant genisligi ve $P_{\text{peak}}$ en yuksek hesaplama kapasitesidir. Fuzyon islemi ara tensor yazma/okuma adimlarini azalttigi icin etkili $I$ degeri artar. Dolayisiyla algoritma, bellek-sinirli bolgeden hesaplama-sinirli bolgeye yaklasir; bu da ayni asimptotik sinifa ragmen daha dusuk calisma suresi ile gozlenir.

\section{Olceklenebilirlik Sonuclari}
Farkli model buyuklukleri (ornegin ResNet18, ResNet50, ResNet101) icin yapilan testler, sabit carpandaki iyilesmenin model boyutuna gore degistigini gostermektedir. Kucuk modellerde I/O maliyeti baskinken, buyuk modellerde hesaplama maliyeti baskin hale gelebilir.

\section{Sonuc}
Bu calisma, ileri algoritma analizi dersi kapsaminda, donanim seviyesindeki graf optimizasyonlarinin teorik modellemeler ile pratik performans arasindaki boslugu nasil doldurdugunu gostermistir. Fuzyon mekanizmasi asimptotik sinifi degistirmese de, I/O kaynakli gizli sabitleri azaltarak ampirik gecikmeyi anlamli olcude iyilestirmektedir.

\bibliographystyle{plain}
\bibliography{references}

\end{document}
